\scanpage{412}\mspage{151}
\diagram{\begin{tikzcd}[ampersand replacement=\&,column sep=large,execute at end picture={\node[left=1.5em] at ($(topL.west)!.5!(lowL.west)$) {$(*)$};}]
|[alias=topL]|\underline{\mathrm{Ch}}^{\deleted{\unsure{W|\infty|\circ}}}_{\infty}(X)_{/(a,b)} \arrow[r,"{\unsure{\rho|p}}"] \& \underline{\mathrm{Ch}}^{\deleted{W}}_{\infty}(X)_{/(a',b')} \\
|[alias=lowL]|\underline{\mathrm{Ch}}^{\deleted{W}}_{\infty}(X;a,b) \arrow[u,"\sim"] \arrow[r,dashed,"\varphi"] \& \underline{\mathrm{Ch}}^{\deleted{W}}_{\infty}(X,\deleted{\unsure{b|a}}a',b') \arrow[u,"\sim","{\unsure{-}}"']
\end{tikzcd}.}
Ainsi on trouve une flèche en pointillés \inserted{dans $\mathrm{Hot}$}, dont
il faut prouver que c'est une \inserted{$\mathrm{Hot}$-}équivalence. Mais
en considérant
\[
u:a\longrightarrow a'\qquad v:b\mathbin{\text{---}}b'
\]
comme des \uline{chemins} dans $X$, et \deleted{\unsure{en|un}}
\[
u^\circ:a'\longleftarrow a
\]
comme un chemin rétrograde de $a'$ à $a$, la
théorie de variance générale (du IX) nous
donne bel et bien un \deleted{\unsure{id}} isomorphisme
dans $(\mathrm{Hot})$ :
\[
\underline{\mathrm{Ch}}_{\infty}(X;a,b)\mathbin{\text{---}}\underline{\mathrm{Ch}}_{\infty}(X;a',b')\quad,
\]
on a envie de dire que c'est $\varphi$, i.e. qu'en
l'insérant \unsure{à|:} la place de $\varphi$, le carré ci-dessus
devient commutatif -- Ce qui prouverait ce
qu'on veut). Je pense que si on travaille
avec les \inserted{ind-objets} $\underline{\underline{\mathrm{Ch}}}_{\infty}(X)_{/(a,b)}$, $\underline{\underline{\mathrm{Ch}}}_{\infty}(X;a,b)$ etc.,
pour lesquels on a des accouplements
\inserted{qui s'}\unsure{explicitent|explicitement}, \unsure{visiblement|strictement} associatifs,
et donnant naissance aux $\varphi$, cela devrait
être facile (quoique un \unsure{peu} fastidieux).
Cela \deleted{\unsure{pr}} ne prouverait pas pour \unsure{autant}
la commutativité de $(*)\deleted{{}_{\unsure{6|\ast}}}$ dans $\mathrm{Hot}$, mais
% archive sheet number: 411
